\section*{Chapitre \ref{ch:ll:the-order-book-builder} --- Le constructeur de carnet d'ordres}

\begin{solution}{exo:ll:the-order-book-builder:1}
La base vaut $50.00 - 2\,048 \times 0.01 = 29.52$~; 49.73 tombe à l'indice $(49.73 - 29.52)/0.01 = 2\,021$. La moitié centrale couvre
les indices 1\,024 à 3\,071, soit des prix de 39.76 à 60.23~: l'échelle reste en place tant que les deux meilleurs prix restent dans cette plage.
\end{solution}

\begin{solution}{exo:ll:the-order-book-builder:2}
$65\,536 \times 16 \times 2 = 2$~MiB, exactement la taille du cache de deuxième niveau du portable (chapitre~3)~: il ne tient que si rien d'autre
ne s'y trouve. En pratique, ce sont le meilleur cours de chaque côté et les quelques niveaux voisins qui restent dans le cache.
\end{solution}

\begin{solution}{exo:ll:the-order-book-builder:3}
La modification retire l'ordre de 500 à 20.10 (son niveau perd 500) et ajoute un nouvel ordre de 300 à 20.11 sous la nouvelle
référence, en queue de la file de ce niveau~: il a perdu sa priorité, comme le prévoit la spécification pour la nouvelle référence.
\end{solution}

\begin{solution}{exo:ll:the-order-book-builder:4}
Avant~: la case 5 contient 1\,024, la case 6 2\,048, la case 7 3\,072. L'effacement de 1\,024 vide la case 5~; 2\,048 (case d'origine 5) revient en case 5 et
3\,072 (case d'origine 5) passe en case 6~; la case 7 est vidée. Un marqueur de suppression aurait laissé la case 5 marquée comme supprimée~: chaque recherche de 2\,048 et de 3\,072
devrait sonder cette case jusqu'à la reconstruction de la table, et après une journée d'annulations, les suites de sondes s'allongeraient.
\end{solution}

\begin{solution}{exo:ll:the-order-book-builder:5}
$\sigma = 40 \times 0.02 / 0.01 = 80$ ticks par jour. Avec 1\,024 ticks, $a = 256$~: $(80/256)^2 \approx 0.098$, soit environ un
recentrage tous les dix jours~; avec 4\,096, $a = 1\,024$~: $0.0061$, soit un recentrage tous les 160 jours.
\end{solution}

\begin{solution}{exo:ll:the-order-book-builder:6}
Il retire l'ordre (en diminuant son niveau des 200 qu'il détient) et compte l'événement comme incohérent~: une réduction supérieure
à l'ordre signifie qu'un événement a été perdu ou dupliqué en amont, et la surveillance doit le faire remonter aux compteurs de
trous du \omterm{def:ll:the-feed-handler:handler}{feed handler}.
\end{solution}

\begin{solution}{exo:ll:the-order-book-builder:7}
Conserver, par côté, un tableau de mots de 64 bits, avec un bit par case de l'échelle, positionné à 1 quand un niveau devient non vide
et remis à 0 quand il se vide. Le prochain meilleur cours acheteur sous l'indice $i$ se trouve mot par mot avec count-leading-zeros sur le mot masqué, le prochain meilleur cours vendeur
avec count-trailing-zeros~: quelques instructions pour 64 cases vides, au lieu de 64 lectures. Le recentrage reconstruit les bitmaps avec les
niveaux.
\end{solution}

\begin{solution}{exo:ll:the-order-book-builder:8}
Une journée calme ne dit rien de celle où le cours se déplace~: une échelle de 512 ticks ($a = 128$) se recentre une fois tous les 128 ticks parcourus,
et sur un titre à petit tick, une journée à 30~\% représente des centaines de milliers de ticks, donc des milliers de recentrages
(l'échelle de 256 ticks de la mesure en demandait près de vingt mille). Dimensionner l'échelle pour le plus grand mouvement attendu, et la tester sur ce mouvement.
\end{solution}

\begin{solution}{pb:ll:the-order-book-builder:1}
\begin{enumerate}
\item $100 \times 0.015 / 10^{-4} = 15\,000$ ticks.
\item $(15\,000/1\,024)^2 \approx 215$ par jour avec 4\,096 ticks~; $(15\,000/16\,384)^2 \approx 0.84$ avec 65\,536.
\item Un recentrage par jour lorsque $a = \sigma$~: $4 \times 15\,000 = 60\,000$ ticks.
\item Le coût de l'échelle vient des journées où le cours parcourt une grande distance~; le taux moyen d'une journée ordinaire les masque.
\item 300\,000 ticks~; une tendance régulière provoque $300\,000/1\,024 \approx 293$.
\item 256~: chaque recentrage centre l'échelle sur le mid, qui a déjà dépassé l'edge, si bien que le suivant exige plus de
$a$ ticks de déplacement.
\item Le carnet lui-même est plus large que la moitié centrale d'une échelle de 256 ticks (des ordres jusqu'à 300 ticks du mid), si bien que les meilleurs
prix quittent l'échelle à chaque saut du meilleur cours, et pas seulement quand le mid se déplace~: 19\,744 recentrages.
\item $(12.8 - 4.4)~\text{ms}/(256 - 15) \approx \qty{35}{\micro\second}$ chacun~: un parcours du pool d'ordres et de l'échelle.
\item 65\,536 ticks~: quinze recentrages et environ \qty{4}{\milli\second} pour la journée.
\item Chaque recentrage vide toute l'échelle, et une échelle de millions de ticks ne tient plus dans aucun cache~: la mise à jour médiane
passe d'environ \qty{29}{\nano\second} à environ \qty{135}{\nano\second} pour $2^{20}$ et $2^{22}$ ticks.
\item 2~MiB.
\item Elle est presque plate, d'environ \qty{30}{\nano\second} de 4\,096 à 65\,536 ticks et de 36 à 262\,144~: dans cette plage, la largeur coûte par
les recentrages, et non par les mises à jour.
\item \textbf{Résultat de référence.} Lors d'une journée à 30~\% sur un titre à petit tick (300\,000 ticks), une échelle de 4\,096 ticks se recentre
256 fois et coûte environ \qty{13}{\milli\second}, une échelle de 65\,536 ticks se recentre quinze fois et coûte \qty{4}{\milli\second}, soit le
minimum de la journée~; une échelle de 256 ticks se recentre près de vingt mille fois et coûte \qty{0.7}{\second}.
\item À partir du cours de l'instrument, de son tick et du plus grand mouvement quotidien attendu (un multiple de sa volatilité)~: quatre fois
ce mouvement en ticks, plafonné par la mémoire et le cache que le carnet peut utiliser.
\item Il sort le recentrage du chemin critique~: les mises à jour continuent sur l'ancienne échelle pendant que la nouvelle se remplit, puis le
constructeur bascule, au prix d'une seconde échelle et d'un moyen de rejouer les mises à jour effectuées entre-temps.
\item Parce qu'une échelle large se recentre ou déborde rarement dans sa table creuse, et les tests doivent exercer ces chemins.
\item Sur un instrument à grand tick, le cours parcourt peu de ticks et le carnet est dense près du meilleur cours~: une petite échelle se
recentre rarement, et chaque structure est rapide.
\item Le nombre de recentrages par jour et le temps que chacun prend, avec l'instrument concerné, afin qu'une largeur devenue trop petite
apparaisse avant de coûter cher.
\item Environ \qty{65}{\nano\second} (vecteur) à \qty{72}{\nano\second} (arbre) par événement sur le flux à petit tick, quel que soit le mouvement~: aucun recentrage, mais le double de la
médiane de l'échelle à chaque événement.
\item Des niveaux à temps constant, au prix de la mémoire et des recentrages~; il faut s'en protéger contre un cours qui parcourt
de grandes distances, par sa largeur et par des tests sur la journée qui compte.
\end{enumerate}
\end{solution}

\begin{solution}{iq:ll:the-order-book-builder:1}
Rechercher un ordre par référence, ajouter à un niveau, le réduire ou le retirer, et lire le meilleur niveau et les premiers niveaux~: tout cela en
dizaines de nanosecondes, à chaque événement. Les ordres du pool sont retrouvés par \omterm{def:ll:the-order-book-builder:open}{adressage ouvert}~; les niveaux sont dans une échelle dense autour du meilleur cours, avec un
repli épars~; les invariants sont vérifiés dans les versions de débogage.
\iqlookfor{opérations, leur fréquence et une structure par opération.}
\end{solution}

\begin{solution}{iq:ll:the-order-book-builder:2}
Un arbre pour la généralité (tout prix, tout tick, sans réglage) à coût logarithmique et une allocation par nouveau niveau~; un vecteur trié
lorsque le carnet est peu profond et que les modifications se concentrent au meilleur cours~; un tableau de niveaux lorsque le tick est connu et que le cours reste
dans une fenêtre, avec un recentrage quand ce n'est pas le cas. Mesurer sur les données propres de l'instrument.
\iqlookfor{le régime de tick et la forme du carnet décident.}
\end{solution}

\begin{solution}{iq:ll:the-order-book-builder:3}
\omterm{def:ll:the-order-book-builder:open}{Adressage ouvert} avec sondage linéaire dans des tableaux plats dimensionnés au démarrage au double des ordres actifs, les valeurs étant des cases dans un
pool préalloué. La suppression décale vers l'arrière les entrées suivantes de la suite de sondes dans le trou lorsque leur case d'origine le permet,
de sorte qu'aucun marqueur de suppression n'allonge les recherches suivantes.
\iqlookfor{tableaux plats, un pool et suppression par décalage vers l'arrière.}
\end{solution}

\begin{solution}{iq:ll:the-order-book-builder:4}
Un événement perdu ou dupliqué (un trou que le \omterm{def:ll:the-feed-handler:handler}{feed handler} n'a pas signalé, une référence inconnue), un bogue du constructeur, ou un croisement légitime
(une enchère). Marquer le carnet comme périmé, cesser de le négocier, vérifier les compteurs de trous du \omterm{def:ll:the-feed-handler:handler}{feed handler} et resynchroniser à partir d'un instantané~;
analyser à partir des événements enregistrés.
\iqlookfor{périmé d'abord, puis la cause, puis la reprise.}
\end{solution}

\begin{solution}{iq:ll:the-order-book-builder:5}
Comparer-le événement par événement à des carnets indépendants sur des jeux de test communs, dans tous les langages où il est écrit~; exercer ses cas rares
chemins (échelles étroites, recentrages, instantanés, références inconnues)~; vérifier les invariants après chaque événement~; rejouer les journées les plus chargées et
les plus volatiles enregistrées~; compter les allocations.
\iqlookfor{test différentiel et chemins rares.}
\end{solution}

\begin{solution}{iq:ll:the-order-book-builder:6}
Avec une volatilité quotidienne $\sigma$ en ticks et un recentrage à $a$ ticks du centre, un cours suivant un mouvement brownien se recentre $\sigma^2/a^2$
fois par jour, et une tendance de $D$ ticks se recentre $D/a$ fois. Dimensionner l'échelle pour le plus grand mouvement attendu, recentrer sur place ou
en arrière-plan, conserver un repli épars, et surveiller les recentrages par instrument.
\iqlookfor{l'argument du temps de sortie et une règle de dimensionnement.}
\end{solution}
